\section{Large Language Model (LLM)}

The Python part of this exercise is available in \myhref{https://github.com/LosDelDGIIM/LosDelDGIIM.github.io/blob/main/subjects/Erasmus-DUE/GKI/Exercises/Exercise12.ipynb}{this Jupyter Notebook}.\\

\begin{ejercicio}[Transformer]
    What is the significance of the attention mechanism in the Transformer architecture? How does it address the limitations of recurrent neural networks for sequence-to-sequence tasks?

    The attention mechanism in the Transformer architecture allows the model to weigh the importance of different input tokens when generating each output token. This enables the model to capture long-range dependencies and relationships between tokens in a sequence, which is a limitation of recurrent neural networks (RNNs) that process sequences sequentially and may struggle with long-term dependencies. By using self-attention, Transformers can process all tokens in parallel, leading to improved efficiency and performance in sequence-to-sequence tasks such as machine translation and text summarization.
\end{ejercicio}

\begin{ejercicio}[Sequence Modeling (Seq2Seq)]
    Compare and contrast the approaches of RNNs and Transformers for sequence modeling tasks such as natural language processing. What are the strengths and weaknesses of each?

    RNNs are designed to process sequential data by maintaining a hidden state that captures information from previous time steps. They are effective for tasks where the order of the input data is important, such as language modeling and time series prediction. However, RNNs can suffer from vanishing or exploding gradients, making it difficult to learn long-term dependencies in sequences.
\end{ejercicio}

\begin{ejercicio}
    What is a token, and what happens during the tokenization and normalization of a text?

    A token is a unit of text, such as a word, subword, or character, that is used as input for natural language processing models. Tokenization is the process of breaking down a text into these individual tokens, while normalization involves standardizing the tokens by converting them to a consistent format (e.g., lowercasing, removing punctuation, etc.). This preprocessing step is essential for preparing text data for analysis and modeling, as it helps reduce variability and improve the model's ability to learn patterns in the data.
\end{ejercicio}

\begin{ejercicio}
    Name three possible stopwords.

    Three possible stopwords are ``the'', ``is'', and ``and''. Stopwords are common words that are often filtered out during text preprocessing because they carry little semantic meaning and can introduce noise into the analysis.
\end{ejercicio}

\begin{ejercicio}
    What is the difference between lemmatization and stemming? Explain using the example of ``were''.

    Lemmatization is the process of reducing a word to its base or dictionary form (lemma), taking into account its meaning and context. For example, the lemma of ``were'' is ``be''. Stemming, on the other hand, is a more crude approach that removes suffixes from words to reduce them to a common root form, often resulting in non-dictionary words. For example, stemming ``were'' might yield ``were''.
\end{ejercicio}

\begin{ejercicio}
    How is a Transformer structured?

    A Transformer is structured with an encoder-decoder architecture. The encoder receives the input sequence and, after embedding and positional encoding, processes it through multiple layers of multi-head attention and feed-forward neural networks.
    
    The decoder takes the encoder's output and the target sequence (shifted right) as input, applying similar layers of multi-head attention and feed-forward networks. This architecture enables the Transformer to effectively model relationships between tokens in a sequence without relying on sequential processing.
\end{ejercicio}

\begin{ejercicio}
    How does TF-IDF work?

    TF-IDF (Term Frequency-Inverse Document Frequency) is a statistical measure used to evaluate the importance of a word in a document relative to a collection of documents (corpus). It combines two components:
    \begin{itemize}
        \item Term Frequency (TF): Measures how frequently a term appears in a document, normalized by the total number of terms in that document.
        \item Inverse Document Frequency (IDF): Measures how important a term is by considering how many documents in the corpus contain that term. It is calculated as the logarithm of the total number of documents divided by the number of documents containing the term.
    \end{itemize}
    The TF-IDF score is obtained by multiplying the TF and IDF values, allowing for the identification of terms that are both frequent in a specific document and rare across the corpus, thus highlighting their significance.
\end{ejercicio}

\begin{ejercicio}
    What does a typical embedder do?

    A typical embedder transforms discrete tokens (e.g., words or subwords) into continuous vector representations (embeddings) in a high-dimensional space. These embeddings capture semantic relationships between tokens, allowing models to understand and process natural language more effectively. Words with similar meanings are represented by vectors that are close together in the embedding space, enabling models to perform tasks such as similarity measurement, clustering, and classification.
\end{ejercicio}

\begin{ejercicio}
    What factors influence the output of an LLM?

    The output of a Large Language Model (LLM) is influenced by several factors, including:
    \begin{itemize}
        \item The model architecture and size (number of parameters).
        \item The quality and diversity of the training data.
        \item The input prompt or context provided to the model.
        \item Fine-tuning or domain-specific training that can adapt the model to specific tasks or knowledge areas.
    \end{itemize}
\end{ejercicio}

\begin{ejercicio}
    What is context size?

    Context size refers to the amount of input data (e.g., tokens or words) that a model can consider at once when generating predictions or outputs. In the case of LLMs, context size determines how much information from the input sequence can be used to inform the model's understanding and response. A larger context size allows the model to capture more dependencies and relationships within the input, leading to more coherent and contextually relevant outputs.
\end{ejercicio}

\begin{ejercicio}
    What does SLM stand for, and how many parameters does it have?

    SLM stands for Small Language Model. The number of parameters in an SLM can vary depending on the specific architecture and design choices, but it typically has fewer parameters than larger models like GPT-3 or BERT.
\end{ejercicio}

\begin{ejercicio}
    What is the difference between explainability and interpretability?

    % // TODO: Hacer
\end{ejercicio}

\begin{ejercicio}
    Name three explainability methods.
\end{ejercicio}